\section{Complete expansions in two asymptotic scales}
\label{sec:expansions}

The positive measure resolves global singularities and signs. We now analyze
its endpoint with enough precision to recover every inverse-logarithmic
coefficient and every fixed algebraic order. The first expansion keeps only
powers of $1/\log n$; the second retains these logarithmic effects inside
convergent integral functions before expanding in $1/n$.

We use $J$, $h$, and $\ell$ from
Section~\ref{sec:spectral-notes}. The boundary value is
$J_+(t)=R(t)+i\pi h(t)$. All truncation indices and all parameters
$a>0$ below are fixed while $n$ or $L$ tends to infinity.

\subsection{The exact analytic endpoint correction}

The stronger identity is
\begin{equation}\label{exp:exactP}
 \frac{J(e^{-s}+i0)}{h(e^{-s})}
 =\log s+\gE+P(s)+i\pi,\qquad 0<s<2\pi,
\end{equation}
where the real-analytic function $P$ has the convergent expansion
\begin{equation}\label{exp:P}
 P(s)=\frac{s^2}{4\sinh^2(s/2)}-1
       +\sum_{k=0}^{\infty}\zeta'(-1-k)\frac{s^{k+2}}{k!}.
\end{equation}
The assertion for the whole interval $(0,2\pi)$ means continuation of the
convergent Taylor expression there. Its local use near zero is all that is
needed below. In particular,
\[
 P(s)=p_2s^2+p_3s^3+p_4s^4+O(s^5),
 \quad
 p_2=\zeta'(-1)-\frac1{12},\quad
 p_3=\zeta'(-2),\quad
 p_4=\frac12\zeta'(-3)+\frac1{240}.
\]

Here is a derivation fixing the signs. Define
\[
 \mathcal B(z)=\sum_{m\geq1}m\log m\,z^m
             =-\left.\partial_\nu\operatorname{Li}_\nu(z)\right|_{\nu=-1}.
\]
The coefficients of $A(z)=\sum_{j\geq1}a_jz^j$ give
\[
 A(z)=(z^{-1}-1)\mathcal B(z)-\frac z{1-z},
 \qquad
 (1-z)(1-A(z))=1-\frac{(1-z)^2}{z}\mathcal B(z).
\]
Differentiate the classical convergent polylogarithm expansion \cite[25.12.12]{DLMF}
\[
 \operatorname{Li}_\nu(e^{-u})
 =\Gamma(1-\nu)u^{\nu-1}
  +\sum_{k\geq0}\zeta(\nu-k)\frac{(-u)^k}{k!},\qquad |u|<2\pi,
\]
at $\nu=-1$, using the principal logarithm away from the negative real axis.
It gives
\[
 \mathcal B(e^{-u})
 =\frac{1-\gE-\log u}{u^2}
  -\sum_{k\geq0}\zeta'(-1-k)\frac{(-u)^k}{k!}.
\]
For $w=e^u$ one has
$J(w)=w(1-w^{-1})(1-A(w^{-1}))$.
Substitute the last display and then let $u\to-s+i0$.
Since $\log(-s+i0)=\log s+i\pi$ and
\[
 h(e^{-s})=e^{-s}\frac{4\sinh^2(s/2)}{s^2},
\]
the result is \eqref{exp:exactP}--\eqref{exp:P}.
The lower boundary has $-i\pi$. The positive sign on the upper boundary also
agrees with the direct Sokhotski--Plemelj formula for $J$.

\subsection{A Laplace form and the correct scale of the first correction}

Let $L=\log n$ and, for $a>0$, define
\begin{equation}\label{exp:Fa}
 F_a(L)=\int_0^\infty
     \frac{y^{a-1}e^{-y}}
      {(L-\log y-\gE)^2+\pi^2}\,dy.
\end{equation}
These functions are positive, finite, and differentiable to every order on
the real axis. The denominator is bounded below by $\pi^2$.
In \eqref{eq:spec-exact-remainder} the change of variables $t=e^{-s}$ is especially useful:
\[
 E_n=\int_0^\infty e^{-ns}\ell(e^{-s})\,ds.
\]
For sufficiently small positive $s$, \eqref{exp:exactP} yields the exact
local density expression
\begin{equation}\label{exp:laplacedensity}
 \ell(e^{-s})
 =\frac{s^2}{1-e^{-s}}
     \frac1{(\log s+\gE+P(s))^2+\pi^2}.
\end{equation}
The part of the original moment integral with $t\leq e^{-\delta}$ is
$O(e^{-(n-1)\delta})$, since $\ell$ is integrable. Thus only a fixed
neighbourhood of $s=0$ enters any of the algebraic orders below.

The expansion $P(s)=O(s^2)$ in \eqref{exp:exactP} gives
\[
 \ell(e^{-s})
 =\frac{s+s^2/2}{(\log s+\gE)^2+\pi^2}
  +O\left(\frac{s^3}{(\log s+\gE)^2+\pi^2}\right).
\]
Consequently
\begin{equation}\label{exp:twosectors}
 E_n=\frac{F_2(L)}{n^2}+\frac{F_3(L)}{2n^3}
           +O\left(\frac1{n^4L^2}\right).
\end{equation}
In particular the first relative correction in powers of $1/n$ is positive.

For comparison, the direct $t$-variable model satisfies the weaker but useful
bound
\[
 E_n=\int_0^1\frac{t^{n-1}(1-t)}
       {(\log(1-t)+\gE)^2+\pi^2}\,dt
         +O\left(\frac1{n^3\log^2n}\right).
\]
Its error has a genuine $n^{-3}\log^{-2}n$ contribution; changing to the
Laplace variable incorporates the corresponding first correction cleanly.

\subsection{Every fixed algebraic order, retaining the logarithmic dependence}

Put
\[
 f(V)=\frac1{V^2+\pi^2},\qquad b(s)=\frac{s}{1-e^{-s}},
 \qquad
 b(s)f(V+P(s))=\sum_{r\geq0}Q_r(V)s^r.
\]
The last identity defines $Q_r$ by Taylor expansion in $s$ while $V$ is held
fixed. Each $Q_r$ is a finite linear combination of derivatives of $f$, with
constants obtained from the Bernoulli coefficients of $b$ and the Taylor
coefficients of $P$. Explicitly,
\begin{align*}
 Q_0&=f,&Q_1&=\tfrac12f,\\
 Q_2&=\tfrac1{12}f+p_2f',&
 Q_3&=(p_2/2+p_3)f',\\
 Q_4&=-\tfrac1{720}f+(p_2/12+p_3/2+p_4)f'
                    +\tfrac12p_2^2 f''.
\end{align*}
Define the coefficient functions
\[
 \Psi_r(L)=\int_0^\infty
   y^{r+1}e^{-y}Q_r(\log y-L+\gE)\,dy.
\]
Then, for every fixed integer $R\geq0$,
\begin{equation}\label{exp:powerexpansion}
 E_n=\sum_{r=0}^{R}\frac{\Psi_r(\log n)}{n^{r+2}}
       +O_R\left(\frac1{n^{R+3}\log^2n}\right).
\end{equation}

Here is a uniform remainder proof. Choose a small fixed circle
$|\xi|=r_0<2\pi$ on which $b$ and $P$ are analytic, and let
$M=\sup_{|\xi|\leq r_0}|P(\xi)|$.
For real $|V|>2(M+\pi+1)$, the function
$b(\xi)f(V+P(\xi))$ is analytic in that disk and is bounded by $C/V^2$.
Cauchy's estimate for its Taylor remainder, for $0<s<r_0/2$, gives
\[
 b(s)f(V+P(s))-\sum_{r=0}^{R}Q_r(V)s^r
       =O_R(s^{R+1}/V^2).
\]
Choose $\delta$ small enough that this applies to
$V=\log s+\gE$ for $0<s<\delta$.
Multiplication by $s$ in \eqref{exp:laplacedensity} makes the local density
error $O_R(s^{R+2}/\log^2s)$.
For each fixed $m\geq0$,
\[
 \int_0^\delta e^{-ns}\frac{s^m}{1+\log^2s}\,ds
       =O_m(n^{-m-1}\log^{-2}n).
\]
Indeed, on $s\leq n^{-1/2}$ the denominator is at least a constant times
$\log^2n$, and the remaining part is exponentially small in $\sqrt n$.
This proves the remainder in \eqref{exp:powerexpansion}.
Finally, every $Q_r$ is bounded on the real axis and decays as $O_r(V^{-2})$,
so extending the coefficient integrals from $y<n\delta$ to $y<\infty$
adds only exponentially small errors.

With the convention that a prime on $F_a$ means differentiation in $L$,
the first four power sectors are therefore
\begin{align}\label{exp:fourpowers}
 E_n={}&\frac{F_2(L)}{n^2}+\frac{F_3(L)}{2n^3}
 +\frac1{n^4}\left\{\frac{F_4(L)}{12}-p_2F_4'(L)\right\}\notag\\
 &\quad-\frac{p_2/2+p_3}{n^5}F_5'(L)
       +O\left(\frac1{n^6L^2}\right).
\end{align}
The minus signs in the derivative terms are forced by
$V=\log y-L+\gE$; $f$ is even but $f'$ is odd.
The $n^{-6}$ coefficient, if wanted, is
\[
 -\frac{F_6}{720}-(p_2/12+p_3/2+p_4)F_6'
                       +\frac{p_2^2}{2}F_6''.
\]
It can be included with an $O(n^{-7}L^{-2})$ remainder by taking $R=4$.

\paragraph{Essential scope distinction.}
In \eqref{exp:powerexpansion}, the functions $\Psi_r$ (or the $F_a$ in
\eqref{exp:fourpowers}) remain intact. Truncating the inverse-logarithmic
expansion of the leading sector at any fixed order produces an error larger
than every later fixed power of $1/n$. A finite rectangular double sum does
not inherit the remainder in \eqref{exp:powerexpansion}.

\subsection{Every fixed inverse-logarithmic order}

For $k\geq0$ put
\[
 q_k(U)=\frac{(U+i\pi)^{k+1}-(U-i\pi)^{k+1}}{2i\pi}
       =\sum_{j=0}^{\lfloor k/2\rfloor}
          (-1)^j\binom{k+1}{2j+1}\pi^{2j}U^{k-2j}.
\]
For every fixed $a>0$ and $K\geq0$,
\begin{equation}\label{exp:logexpansion}
 F_a(L)=\Gamma(a)\sum_{k=0}^{K}\frac{c_{a,k}}{L^{k+2}}
                +O_{a,K}(L^{-K-3}),\qquad L\to\infty,
\end{equation}
where
\[
 c_{a,k}=\frac1{\Gamma(a)}\int_0^\infty
       y^{a-1}e^{-y}q_k(\log y+\gE)\,dy.
\]
To prove the remainder, split the integral into
$y<e^{-L/2}$, $e^{-L/2}\leq y\leq e^{L/2}$, and $y>e^{L/2}$.
The outer parts, including the tails of the proposed coefficients, are
smaller than every fixed inverse power of $L$.
On the middle part expand
\[
 \frac1{(L-U)^2+\pi^2}
 =\frac1{2i\pi}\left\{\frac1{L-U-i\pi}-\frac1{L-U+i\pi}\right\}.
\]
For sufficiently large $L$, $|U\pm i\pi|\leq3L/4$ there.
The finite geometric expansion through the required order has a remainder
bounded by $C_K L^{-K-3}(1+|U|^{K+2})$.
All logarithmic moments against $y^{a-1}e^{-y}$ are finite, proving
\eqref{exp:logexpansion}.

The coefficient generating function has a particularly simple form:
\begin{equation}\label{exp:coeffgen}
 \sum_{k\geq0}\frac{c_{a,k}}{(k+1)!}z^k
 =\frac{\Gamma(a+z)}{\Gamma(a)}e^{\gE z}
          \frac{\sin\pi z}{\pi z}.
\end{equation}
This follows first for small $z$ by integrating
$\sum q_k(U)z^k/(k+1)!=e^{Uz}\sin(\pi z)/(\pi z)$.
For a positive integer $a$ the reflection formula rewrites it as
\[
 \sum_{k\geq0}\frac{c_{a,k}}{(k+1)!}z^k
 =\left\{\prod_{j=1}^{a-1}(1+z/j)\right\}
                       \frac{e^{\gE z}}{\Gamma(1-z)}.
\]
The right-hand side is entire. This is a coefficient-generation statement;
it does not assert convergence of the asymptotic inverse-log series.

For the leading sector, $a=2$, let $c_k=c_{2,k}$. Then
\begin{equation}\label{exp:maincoeffgen}
 \sum_{k\geq0}\frac{c_k}{(k+1)!}z^k
 =(1+z)\frac{e^{\gE z}}{\Gamma(1-z)}
 =(1+z)\exp\left\{-\sum_{m\geq2}\frac{\zeta(m)}m z^m\right\}.
\end{equation}
The displayed exponential identity holds in $|z|<1$; the preceding gamma
quotient gives its entire continuation. The first five coefficients are
\[
 c_0=1,\quad c_1=2,\quad c_2=-\frac{\pi^2}{2},\quad
 c_3=-2\pi^2-8\zeta(3),\quad
 c_4=\frac{\pi^4}{12}-40\zeta(3).
\]
For a recursive implementation set $b_0=1$, $b_1=0$ and
\[
 b_k=-\frac1k\sum_{m=2}^{k}\zeta(m)b_{k-m}\quad(k\geq2),
 \qquad c_k=(k+1)!(b_k+b_{k-1}),
\]
where $b_{-1}=0$. Thus, for every fixed $K$,
\begin{equation}\label{exp:mainresult}
 g_n=\gamma\rho^{-n}
 -\frac1{n^2}\sum_{k=0}^{K}\frac{c_k}{\log^{k+2}n}
 +O_K\left(\frac1{n^2\log^{K+3}n}\right).
\end{equation}
In particular,
\[
 g_n-\gamma\rho^{-n}
 =-\frac1{n^2\log^2n}
 \left\{1+\frac2{\log n}-\frac{\pi^2}{2\log^2n}
       -\frac{2\pi^2+8\zeta(3)}{\log^3n}
       +O(\log^{-4}n)\right\}.
\]
For the $n^{-3}$ sector in \eqref{exp:twosectors}, the normalized first
coefficients are
\[
 \frac12F_3(L)=\frac1{L^2}
 \left\{1+\frac3L+\frac{3-\pi^2/2}{L^2}
              -\frac{3\pi^2+8\zeta(3)}{L^3}+O(L^{-4})\right\}.
\]
Their generating function is
$(1+z)(1+z/2)e^{\gE z}/\Gamma(1-z)$.

